\begin{english}
\noindent\textenglish{This thesis develops and evaluates a data-driven framework for the estimation and clustering of electromechanical oscillatory modes in power systems using ambient data. The subspace identification method Numerical Algorithms for Subspace State Space System Identification (N4SID) is employed to extract modal frequencies and damping, while seven clustering algorithms are comparatively evaluated in terms of their ability to group the resulting estimates and identify the dominant system modes. The methodology is applied to the IEEE 39-Bus benchmark system, while a complementary disturbance-based analysis using the Matrix Pencil (MP) method is conducted on the Kundur Two-Area and IEEE 39-Bus benchmark systems. The comparative evaluation identifies Density-Based Spatial Clustering of Applications with Noise (DBSCAN) as the most accurate algorithm in terms of agreement with the reference modes obtained from PowerFactory, whereas Hierarchical DBSCAN (HDBSCAN) provides the best balance among clustering quality, data coverage, and noise rejection. Finally, the methodology is examined both in a conventional power system and in a power system incorporating wind generation units based on Doubly-Fed Induction Generators (DFIGs). The results reveal changes in modal frequencies and in the participation patterns of the synchronous generators, without materially altering the relative performance of the clustering methods under consideration.}

\paragraph{\textenglish{Keywords:}} \textenglish{\textit{power systems, electromechanical oscillatory modes, system identification, clustering, Matrix Pencil, N4SID, DFIG, Kundur Two-Area System, IEEE 39-Bus System, ringdown analysis, ambient analysis, inter-area, intra-area}}
\end{english}
